% ===========================================================================
%  Chapter 11 notes — Simulation
%  Ross, Introduction to Probability Models (10th ed.)
%
%  DIVISION OF LABOUR (see AI_INSTRUCTIONS.md §7):
%    The assistant owns the scaffold: preamble, sectioning, and the faithful
%    transcription of definitions, theorem statements, and examples.
%    The user owns everything between the NOTE / SOLUTION markers.
% ===========================================================================
\documentclass[11pt]{article}
\usepackage{coursemacros}

\title{Chapter 11 --- Simulation}
\author{Mikey Ferguson}
\date{\today}

\begin{document}
\maketitle
\tableofcontents
\bigskip

% ---------------------------------------------------------------------------
\section{Overview}
% ===== NOTE 11.overview =====
% TODO(mferguson): what this chapter is for, in your own words.
% ===== END NOTE 11.overview =====

% ---------------------------------------------------------------------------
\section{Definitions and setup}
% Assistant: transcribe definitions and the model setup here, in order.

% ---------------------------------------------------------------------------
\section{Results}
% Assistant: transcribe theorem/proposition STATEMENTS here using the
% booktheorem environment. Derivations are the user's work.

% ---------------------------------------------------------------------------
\section{Lecture module notes}
% The professor's module slides for this chapter live in this chapter's
% lectures/ folder. Anything from lecture that is not in Ross goes here.
% ===== NOTE 11.lectures =====
% TODO(mferguson): what lecture added or emphasised beyond the text.
% ===== END NOTE 11.lectures =====

% ---------------------------------------------------------------------------
\section{Worked examples}
% ===== NOTE 11.examples =====
% TODO(mferguson): the small examples you worked by hand while reading.
% ===== END NOTE 11.examples =====

% ---------------------------------------------------------------------------
\section{Loose ends}
% ===== NOTE 11.questions =====
% TODO(mferguson): what is still unclear.
% ===== END NOTE 11.questions =====

\end{document}
